\section{Appendix}
\label{sec:appendix}

\subsection{Proof of Theorem 1}

\begin{theorem}[Pruning via Abstract Expressions]
\label{append:thm:pruning}
For an input \graph $G_0$, and a \graph $G$ equivalent to $G_0$,
% Let $g$ be the abstract expression of the input program and $G$ a \graph that is equivalent to the input program.
if $\Aeq \models E(G_0) = E(G)$ then $G$ will be generated by \Cref{algo:generate} and will not be pruned away.
\end{theorem}

\begin{proof}
By \Cref{tab:expression,tab:axioms}, we show that for any operator $\er{op}$, if $Y = \er{op}(X_1,\ldots,X_n)$ then 
$\Asub \models \subexpr(\expr(X_i), \expr(Y))$ (for $1 \leq i \leq n)$.
That is, an input to an operator is always an abstract subexpression of its output.
By the reflexivity and transitivity axioms included in $\Asub$,
it follows for any $G'$ that is a prefix of $G$, $\Asub \models \subexpr(\expr(G'),\expr(G))$.
Together with the assumption that $\Aeq \models E(G_0) = E(G)$,
it follows that $\Aeq \cup \Asub \models \subexpr(\expr(G'),\expr(G_0))$.
Thus, any prefix of $G$ will not be pruned, and $G$ will be generated by \sys.
\end{proof}
